\chapter{Eyelid Bump (Stye and Chalazion)}

\section*{Definition}
A localized, tender, erythematous nodule on the eyelid margin caused by acute bacterial infection of eyelash follicle (external hordeolum/style) or meibomian gland (internal hordeolum); a chalazion is a chronic, painless, inflammatory granuloma of the meibomian gland.

\section*{When to See a Doctor}
See your doctor if:
\begin{itemize}
  \item Recurrent styes, lesion not resolving within 2 weeks, vision changes, eyelid swelling spreading to entire eye or face, or fever
\end{itemize}

\section*{How It Develops}
External styes involve infection of the Zeis or Moll glands (eyelash follicles) by Staphylococcus aureus, causing acute suppurative inflammation. Internal styes infect meibomian glands with similar pathophysiology. Chalazion results from obstruction of the meibomian gland duct without acute infection, leading to lipogranulomatous inflammation with lipid-laden macrophages, giant cells, and a firm, nontender nodule.

\section*{Causes}
\begin{itemize}
  \item Staphylococcus aureus infection (most common cause of stye)
  \item Poor eyelid hygiene
  \item Blepharitis (chronic eyelid inflammation predisposing to styes)
  \item Contact lens use (improper hygiene)
  \item Rosacea (associated with meibomian gland dysfunction)
  \item Seborrheic dermatitis
  \item Hormonal changes
  \item Stress and fatigue
  \item Use of expired or contaminated eye makeup
  \item Duct obstruction from thickened meibomian secretions (chalazion)
\end{itemize}

\section*{Related Symptoms}
\begin{itemize}
  \item Eyelid swelling
  \item Eye redness
  \item Eye discharge
  \item Eye pain
  \item Blepharitis
\end{itemize}

\section*{Medical Evaluation}
\begin{itemize}
  \item Your doctor will examine the eyelid bump to differentiate a stye (tender, erythematous) from a chalazion (firm, painless).
  \item Eversion of the eyelid may be performed to assess the inner conjunctival surface.
  \item Visual acuity is checked to ensure the bump is not affecting vision.
  \item Slit-lamp examination evaluates the eyelid margin and meibomian glands.
  \item Referral to an ophthalmologist is indicated for recurrent lesions or those not responding to conservative treatment.
\end{itemize}

\section*{Treatment Options}
\begin{itemize}
  \item Warm compresses and gentle massage are first-line for both styes and chalazia.
  \item Topical antibiotics are not routinely needed; an ophthalmic clinician may prescribe them when infection or associated blepharitis warrants treatment.
  \item Oral antibiotics are reserved for selected associated conditions, such as significant rosacea or spreading infection, under clinician guidance.
  \item Incision and curettage for chalazia that persist despite conservative therapy.
  \item Eyelid hygiene and warm compresses to prevent recurrence.
\end{itemize}


\section*{Self-Care and Home Management}
\begin{itemize}
  \item Apply warm compresses to the affected eyelid for 10-15 minutes, 3-4 times daily.
  \item Gently massage the bump with clean fingers after warm compresses to promote drainage.
  \item Maintain good eyelid hygiene by cleaning the lid margins with baby shampoo or lid wipes.
  \item Avoid squeezing or popping the bump, as this can spread infection.
  \item Do not wear eye makeup or contact lenses until the bump resolves.
\end{itemize}

\section*{References}
\begin{thebibliography}{99}
\bibitem{eyelid_bump:lindsley2010}
Lindsley K, Nichols JJ, Dickersin K. Interventions for acute internal hordeolum. \textit{Cochrane Database Syst Rev}. 2010;(9):CD007742.
\bibitem{eyelid_bump:gilchrist2008}
Gilchrist H, Lee GA. Management of chalazia and hordeola. \textit{Aust Fam Physician}. 2008;37(10):823--825.
\bibitem{eyelid_bump:kanski2015ref}
Kanski JJ, Bowling B. Eyelid disorders. In: \textit{Clinical Ophthalmology}. 8th ed. Elsevier; 2016.
\bibitem{eyelid_bump:american2017ref}
American Academy of Ophthalmology. \textit{Preferred Practice Pattern: Blepharitis}. AAO; 2013.
\bibitem{eyelid_bump:bartley1998}
Bartley GB. The differential diagnosis and classification of eyelid tumors. \textit{Int Ophthalmol Clin}. 1998;38(2):1--13.
\bibitem{eyelid_bump:lederman1999}
Lederman C, Miller M. Hordeola and chalazia. \textit{Pediatr Rev}. 1999;20(8):283--284.
\end{thebibliography}
